% mooie C++
% http://archive.cs.uu.nl/mirror/CTAN/macros/latex/contrib/relsize/relsize-doc.pdf
\usepackage{relsize}
\iftoggle{charter}{
    \newcommand{\Cpp}{C\nolinebreak\raisebox{.2ex}[0pt][0pt]{\smaller{+\nolinebreak+}}}
}{
\newcommand{\Cpp}{C\nolinebreak\raisebox{.3ex}[0pt][0pt]{\smaller{+\nolinebreak+}}}
}

% http://archive.cs.uu.nl/mirror/CTAN/macros/latex/required/graphics/grfguide.pdf
\usepackage{color}
% define colors to use in listings
\definecolor{lstblue}{RGB}{42,0,255}
\definecolor{lstpurple}{RGB}{127,0,116}
\definecolor{lstturquoisegreen}{RGB}{63,127,95}

\newcommand{\lstcppstring}{\color{lstblue}}
\newcommand{\lstcppkeyword}{\color{lstpurple}\bfseries}
\newcommand{\lstcppincludefile}{\color{lstblue}}
\newcommand{\lstcppcomment}{\color{lstturquoisegreen}}

% http://archive.cs.uu.nl/mirror/CTAN/macros/latex/contrib/listings/listings.pdf
\usepackage{listings}
\lstset{
    inputencoding={utf8}, 
    literate=
        {ä}{{\"a}}1 
        {ë}{{\"e}}1 
        {ï}{{\"i}}1 
        {ö}{{\"o}}1
        {ü}{{\"u}}1
        {á}{{\'a}}1
        {é}{{\'e}}1
        {í}{{\'i}}1
        {ó}{{\'o}}1
        {ú}{{\'u}}1
        {à}{{\`a}}1
        {è}{{\`e}}1
        {ì}{{\`i}}1
        {ò}{{\`o}}1
        {ù}{{\`u}}1
        {â}{{\^a}}1
        {ê}{{\^e}}1
        {î}{{\^i}}1
        {ô}{{\^o}}1
        {û}{{\^u}}1,
    language=C++,
    morekeywords={alignas,alignof,char16_t,char32_t,concept,constexpr,decltype,noexcept,nullptr,requires,static_assert,thread_local,override,final},
    morekeywords=[2]{<cstdlib>,<csignal>,<csetjmp>,<cstdarg>,<typeinfo>,<typeindex>,<type_traits>,<bitset>,<functional>,<utility>,<ctime>,<chrono>,<cstddef>,<initializer_list>,<tuple>,<any>,<optional>,<variant>,<compare>,<version>,<new>,<memory>,<scoped_allocator>,<memory_resource>,<climits>,<cfloat>,<cstdint>,<cinttypes>,<limits>,<exception>,<stdexcept>,<cassert>,<system_error>,<cerrno>,<contract>,<cctype>,<cwctype>,<cstring>,<cwchar>,<cuchar>,<string>,<string_view>,<charconv>,<array>,<vector>,<deque>,<list>,<forward_list>,<set>,<map>,<unordered_set>,<unordered_map>,<stack>,<queue>,<span>,<iterator>,<ranges>,<algorithm>,<execution>,<cmath>,<complex>,<valarray>,<random>,<numeric>,<ratio>,<cfenv>,<bit>,<iosfwd>,<ios>,<istream>,<ostream>,<iostream>,<fstream>,<sstream>,<syncstream>,<strstream>,<iomanip>,<streambuf>,<cstdio>,<locale>,<clocale>,<codecvt>,<regex>,<atomic>,<thread>,<mutex>,<shared_mutex>,<future>,<condition_variable>,<filesystem>},
    basicstyle=\ttfamily\small,
    tabsize=4,
    stringstyle=\lstcppstring,
    keywordstyle=\lstcppkeyword,
    alsoletter={<>.},
	keywordstyle=[2]\lstcppincludefile,
    identifierstyle=,
    commentstyle=\lstcppcomment,
    showstringspaces=false,
    upquote=true,
    numbers=none,
    aboveskip=.1\baselineskip,
    belowskip=0pt,
    abovecaptionskip=.3\baselineskip,
    belowcaptionskip=-.2\baselineskip,
    rangebeginprefix=//>,
    rangeendprefix=//<,
    includerangemarker=false,
    prebreak={\raisebox{0ex}[0ex][0ex]{\textcolor{black}{\ensuremath{\hookleftarrow}}}}, 
    postbreak={\raisebox{0ex}[0ex][0ex]{\textcolor{black}{\ensuremath{\space\hookrightarrow\space}}}}, 
    breaklines=true, 
    breakatwhitespace=true, 
    breakindent=0.1em,
    numberbychapter=true,
    captionpos=b,
}

% \intextcodefont is nodig om \small te verwijderen
\newcommand{\intextcodefont}{\ttfamily}
\newcommand{\code}[2][]{%
	\lstinline[
		basicstyle=\intextcodefont, 
		prebreak={}, 
		postbreak={}, 
		breakindent=0pt,
        #1
	]@#2@%
}

% met \hcode kun je ercoor zorgen dat lange identifiers afgebroken kunnen worden. Mogelijke afbreekpunten met een - markeren.
% TODO soms valt - weg ?
\newcommand{\hcode}[2][]{%
    \lstinline[
        basicstyle=\ttfamily, 
        literate={-}{}{0\discretionary{-}{}{}},
        prebreak={}, 
        postbreak={}, 
        breakindent=0pt, 
        breakatwhitespace=false,
        #1
    ]@#2@%
}

% handige afkorting |bla| voor \code{bla} maar pas op: werkt niet in footnotes, captions etc
\lstMakeShortInline[
	basicstyle=\intextcodefont, 
	prebreak={}, 
	postbreak={}, 
	breakindent=0pt
]{|}

% gebruik voor opmaak van warnings en error meldingen van de C++ compiler
\newcommand{\error}{\texttt}
% gebruik voor opmaak van console output
% http://ftp.snt.utwente.nl/pub/software/tex/macros/latex/contrib/fancyvrb/fancyvrb.pdf
% nodig om underline voor input te kunnen gebruiken: \cinput{...}
\usepackage{fancyvrb}
\newcommand{\cinput}[1]{\color{green}{\underline{#1}}}
\newenvironment{coutput}{%
	\Verbatim[commandchars=\\\{\}]%
	}{%
	\endVerbatim
}

\newcommand{\proglink}[1]{%
    \href{\progsurl/#1}{#1}%
}
